% ======= LaTeX Setup ======= %
\documentclass[11pt]{article}
\usepackage[margin=1in]{geometry}
\usepackage{times}
\usepackage{indentfirst}
\usepackage{booktabs}
\usepackage{tabularx}
\usepackage{float}
\usepackage{amsmath}
\newcolumntype{L}{>{\raggedright\arraybackslash}X}
\linespread{1.15}
\setcounter{tocdepth}{2}
\usepackage[hidelinks]{hyperref}

\title{On-Device Classification of Neural Action Potentials for\\
Wireless Transmission over Bluetooth Low Energy}
\author{Fredrik Willumsen Haug \\[0.6em]
{\normalsize Project Advisor: Dr.\ Saehyuck Oh} \\[0.2em]
{\normalsize SEAS Faculty Thesis Reader: Professor Gage Hills}}
\date{}

\begin{document}

\maketitle

\tableofcontents
\bigskip

% ======= 0. Executive Summary ======= %
\setcounter{section}{-1}
\section{Executive Summary}
This thesis addresses a bandwidth bottleneck in wireless neural recording. Bluetooth Low
Energy (BLE) 6.0, the leading low-power wireless protocol, has a 4~Mbps theoretical ceiling
but reaches only $\sim$1~Mbps in practice. Extracellular action potentials contain
frequencies from around 300~Hz to 6~kHz and are sampled at 30~kHz as the standard in the
field. At 16-bit resolution, each channel produces 480~kbps. BLE can therefore stream raw
data from only two channels which is far below modern high-density probes with hundreds to
thousands of channels. Instead of transmitting raw data for offline spike sorting, this
project detects and classifies spikes on-device, transmitting only a unit ID and timestamp
per spike. This cuts the per-spike data volume by orders of magnitude, enabling far more
channels to be streamed wirelessly on a commercial BLE system-on-chip which is a capability
not previously demonstrated on commercial hardware.

% ======= 1. Problem Identification ======= %
\section{Problem Identification}

\subsection{Background and Motivation}
Whether the substance of thought is reducible to physical processes remains an open
question. Nevertheless, evidence undisputedly suggests that the organized interplay between
neurons in the central nervous system gives rise to human sensory and cognitive
capacity~[1]. The first recording of what is now known as an action potential or
spike---the voltage across an electrically excitable cell's membrane as a function of
time---was published in 1868 by Julius Bernstein~[2]. Derived from frog tissue using a
differential rheotome, an apparatus he invented, Bernstein's measurement provided pivotal
evidence for theories regarding neural dynamics propounded by physicists and physiologists
preceding him~[2]. The interplay between neurons can, in turn, be studied by measuring the
action potentials of individual neurons within and across anatomically defined brain
regions~[3].

A desire to measure action potentials at increased scales has motivated the development of
high-density silicon recording probes. The Neuropixels 1.0 probe, introduced in 2017, is an
intracranial high-density probe for extracellular electrophysiology recordings in rodents
equipped with 384 channels for simultaneous recording~[3]. Its successor, the Neuropixels
2.0, was presented in 2021 and extends recording capacity to as many as 5120
channels~[4]. Such systems have enabled groundbreaking discoveries in systems neuroscience,
and the Neuropixels 3.0 is currently being developed~[5, 6]. Higher channel counts,
however, impose a proportionally higher bandwidth requirement.

The industry-standard analog front end for sampling raw extracellular voltage uses a 16-bit
analog-to-digital converter (ADC) typically configured with a 30~kHz sampling rate~[7,
8]. At these settings, a full Neuropixels 2.0 quad-base recording produces a data rate of
$\sim$2.46~Gbps~[A1]. Bluetooth Low Energy 6.0 (BLE) is a state-of-the-art consumer
wireless communication protocol with a power budget suitable for an implantable
device~[9, 10]. In contrast to the data rate generated by state-of-the-art recording
probes, BLE has a theoretical maximum data rate of 4~Mbps, just $\sim$0.16\% of the raw
data rate generated by a Neuropixels 2.0 quad-base recording~[A2]. Fortunately, many
questions in neuroscientific inquiry require only neural spike signatures, constituting a
small fraction of the continuous raw signal~[11, 12].

\subsection{Stakeholders, Client, and End User of the Solution}
The primary stakeholders are systems neuroscientists who study how populations of neurons
give rise to behavior. Their central experimental paradigm correlates neural activity with
what an animal is doing, which requires the animal to move and behave freely. Present
high-density recording, however, tethers the animal to an acquisition system by a cable
that carries the raw data stream, constraining natural movement, precluding many social and
naturalistic assays, and limiting how many animals can be recorded at once. These
researchers would directly benefit from a recording head stage that transmits the neural
information relevant to their analysis wirelessly, freeing the animal from the tether while
preserving the spike-level detail their experiments depend on. They are therefore both the
end users of the proposed solution and the community whose scientific questions motivate
it.

\subsection{Problem Statement}
The bandwidth of Bluetooth Low Energy 6.0 is too low to support transmission of raw neural
extracellular electrophysiology data sampled at the industry standard from more than two
recording channels at once for offline spike sorting and consequent neuroscientific
analysis~[21].

% ======= 2. Existing Solutions and Previous Work ======= %
\section{Existing Solutions and Previous Work}

\subsection{Existing Solutions}
Existing approaches to the bandwidth problem fall into three groups: transmitting raw or
lightly-processed data for offline sorting, classifying spikes on-device on hardware
unsuited to wireless deployment, and calibrating templates offline to classify on-device in
simulation only. The dominant paradigm transmits raw or lightly processed data and sorts
offline. Headstages with off the shelf components either stream the raw signal and sort
entirely offline~[13] or detect and wavelet-compress spikes on-device yet still cluster
them at the receiver~[14]. In both, the wireless link carries waveforms, not neural
identities. On-device classification exists, but on hardware unsuited to this
goal. Fabricated custom-silicon application-specific integrated circuits (ASICs) reach
85--99\% on-chip accuracy~[15, 16, 17] but are inaccessible parts validated only on
replayed data with no radio, while the closest commodity systems sort on programmable logic
run real-time template matching on an Intan-plus-FPGA headstage in a freely-moving
primate~[18], Wu et al.~[19] run unsupervised OSort on an FPGA yet the former logs to an SD
card and neither integrates a wireless link. Finally, the two-stage scheme of calibrating
templates offline and classifying on-device so that only identities are transmitted was
shown feasible only in simulation~[20]. No prior system, to my knowledge, performs
on-device single-unit classification and wirelessly transmits the resulting spike
identities on a commodity Bluetooth Low Energy SoC.

\subsection{Previous Work}
Work completed so far comprises a single printed circuit board. A custom printed circuit
board for testing the signal path, carrying a bonded Intan RHD2164 analog front end
interfaced to an nRF54LM20B SoC over SPI. The board has been brought up and used to
validate end-to-end single-channel acquisition. A known analog input digitized by the Intan
and relayed over the SoC's Bluetooth Low Energy link. Establishing the acquisition front end
on which the remaining functional requirements will be built.

\subsection{Proposed Solution}
This project proposes a recording head stage that detects spikes and classifies each to a
single unit on-device, transmitting only a compact unit identity and timestamp per spike
rather than the raw waveform. Built entirely on off-the-shelf hardware---an Intan analog
front end paired with a commercial Bluetooth Low Energy system-on-chip---it reduces the
per-spike data volume by orders of magnitude, allowing far more channels to be streamed
wirelessly than raw transmission permits while preserving the spatiotemporal spike
information that neuroscientific analysis depends on.

% ======= 3. Design Requirements and Technical Specifications ======= %
\section{Design Requirements and Technical Specifications}

\subsection{Description}
This section decomposes the problem statement into the behaviors any valid solution must exhibit
(the functional requirements) and the measurable qualities those behaviors must satisfy (the
technical requirements). Furthermore, the specific values for the technical requirements
(technical specifications) are asserted. Requirements and specifications are stated
independently of any particular design; the design approaches that could realize them will
be considered separately in a design section.

Compressing continuous neural data without fidelity loss of spatiotemporal spike
information requires isolating and preserving exactly the neural spike events at the
required resolution. Any solution to the problem statement must therefore
perform five operations: it must \emph{acquire} extracellular voltage from multiple
recording channels, \emph{detect} spike events within the acquired signal, \emph{log} the
time at which each detected spike occurs, \emph{classify} each detected spike and assign it
a unit identity, and \emph{transmit} each unit identity coupled with its timestamp over a
wireless link. These five behaviors constitute the functional requirements of the system
summarized in Table~\ref{tab:funcreq}.

% ======= Table 1 ======= %
\begin{table}[H]
\centering
\caption{Functional requirements of a proposed system}
\label{tab:funcreq}
{\footnotesize
\setlength{\tabcolsep}{4pt}
\renewcommand{\arraystretch}{1.25}
\begin{tabularx}{\textwidth}{@{}l L@{}}
\toprule
Label & Functional requirement \\
\midrule
F1 & \emph{Acquire} extracellular voltage from multiple recording channels \\
F2 & \emph{Detect} spike events within the acquired signal \\
F3 & \emph{Log} the time at which each detected spike occurs \\
F4 & \emph{Classify} each detected spike and assign it a unit identity \\
F5 & \emph{Transmit} each unit identity coupled with its timestamp over a wireless link \\
\bottomrule
\end{tabularx}
}
\end{table}

\subsection{Technical Requirements}
Each functional requirement generally further imposes one or more technical requirements
grouped by the behavior that instantiates a technical requirement. Each technical
requirement follows from one of three sources: a first-principles calculation, a value
established in the prior-art literature, or preliminary testing conducted for this
project. The technical requirements are summarized as T1 through T11 grouped by functional
requirement in Table~\ref{tab:techreq}. Accompanying justification of the derived technical
requirements follows after Table~\ref{tab:techreq}.

% ======= Table 2 ======= %
\begin{table}[H]
\centering
\caption{Technical requirements with specifications grouped by functional requirements}
\label{tab:techreq}
{\footnotesize
\setlength{\tabcolsep}{4pt}
\renewcommand{\arraystretch}{1.25}
\begin{tabularx}{\textwidth}{@{}l L L@{}}
\toprule
Label & Requirement & Specification \\
\midrule
\multicolumn{3}{@{}l}{\textbf{F1 --- Acquire}} \\
\addlinespace[2pt]
T1  & $\geq 3$ channels amplified and digitized & $\geq 3$; target 32 \\
T2  & Sample fast enough for the spike band & $\geq 12.5$~kS/s; design 25~kS/s \\
T3  & Digitization keeps up across channels & $\geq$ channels $\times$ per-channel rate \\
T4  & Resolve spikes above the noise floor & 16-bit; SNR [TBD] \\
T5  & Link carries the digitized stream & $\geq$ channels $\times$ rate $\times$ bits \\
\midrule
\multicolumn{3}{@{}l}{\textbf{F2 --- Detect}} \\
\addlinespace[2pt]
T6  & Spikes separated from the LFP & $\sim$300~Hz \\
T7  & Detect spikes reliably & TPR [TBD]; FPR [TBD] \\
\midrule
\multicolumn{3}{@{}l}{\textbf{F3 --- Log}} \\
\addlinespace[2pt]
T8  & A timestamp is recorded per spike & one per detected spike \\
\midrule
\multicolumn{3}{@{}l}{\textbf{F4 --- Classify}} \\
\addlinespace[2pt]
T9  & Each spike assigned a unit identity & one of $K$ units \\
T10 & Sort better than chance & $> 1/K$ (chance); [TBD] \\
\midrule
\multicolumn{3}{@{}l}{\textbf{F5 --- Transmit}} \\
\addlinespace[2pt]
T11 & Packet carries identity and timestamp & 5~B (1~B ID $+$ 4~B ts) \\
\bottomrule
\end{tabularx}
}
\end{table}

\emph{Acquisition (F1)} imposes T1--T5. As established in the problem statement, Bluetooth
Low Energy cannot carry raw data beyond two channels. Recording from at least three
channels (T1) is therefore the minimum that exceeds this baseline. Extracellular spike
energy occupies approximately 300~Hz--6~kHz~[22], so by the Nyquist--Shannon sampling
theorem faithful capture requires a per-channel rate of at least $\sim$12~kS/s~[A3]; the
30~kS/s field standard~[7, 8] (T2) adds $\sim$3$\times$ oversampling margin accounting for
non-perfections of the bandpass filter. The digitization must keep pace with every channel
at once, so the aggregate conversion throughput must be at least the product of channel
count and per-channel rate (T3). Spike amplitudes of tens to hundreds of microvolts must be
resolved above a microvolt-scale noise floor, for which the industry-standard 16-bit
converter~[7, 8] provides sufficient dynamic range (T4). Finally, the digitized stream must
reach the processor without loss, so the link throughput must meet or exceed the aggregate
data rate (T5); the same investigation identified link throughput as the binding constraint
on channel count, as the serial interfaces of commercial BLE SoCs span roughly 8--32~Mbps
depending on the instance used~[23].

\emph{Detection (F2)} imposes T6--T7. Spikes are superimposed on the local field potential,
which occupies frequencies below $\sim$300~Hz~[22]; a high-pass corner near 300~Hz
therefore isolates the spike band (T6). Detection must then recover a useful fraction of
true spikes without over-triggering on noise, benchmarked against the established Kilosort
sorter as ground truth~[24] (T7); fabricated on-chip sorters report detection accuracies
above 97\%~[15, 16], setting the reference against which the detection target is calibrated.

\emph{Timekeeping (F3)} imposes T8. Because the proposed solution preserves
spatiotemporal spike information, each detected spike must carry a timestamp referenced to
the acquisition clock---one per event being the minimal sufficient encoding (T8).

\emph{Classification (F4)} imposes T9--T10. Each spike is assigned to one of $K$ resolved
units (T9). As a minimum viable criterion, on-device assignment must exceed chance---$1/K$
for $K$ balanced units---relative to offline Kilosort labels~[24] (T10). This floor is
deliberately conservative: comparable on-chip sorters achieve classification accuracies of
roughly 85--99\%~[15, 16, 17], which sets the aspirational reference, while exceeding chance
is the falsifiable bar that proves the pipeline classifies rather than guesses.

\emph{Transmission (F5)} imposes T11. The compression the proposed solution demands is realized by
transmitting, per spike, only a unit identity and a timestamp rather than the raw waveform:
a one-byte identity together with a four-byte timestamp---a 32-bit counter that wraps only
after $\sim$40~hours at 30~kHz---total five bytes per event, against the roughly eighty
bytes a raw forty-sample snippet would require~[A4] (T11).

The three values left as [TBD]---the front-end SNR (T4), the detection sensitivity and
false-positive rate (T7), and the classification-accuracy target (T10)---are quantities
better fixed by direct measurement than by assertion; the corresponding measurement methods
are documented separately.

% ======= Appendix ======= %
\section*{Appendix}
\addcontentsline{toc}{section}{Appendix}

\noindent\textbf{[A1] Raw Neuropixels 2.0 data rate.} The aggregate raw rate is channel
count $\times$ per-channel sample rate $\times$ bit depth:
$5120 \times 30\text{~kS/s} \times 16\text{~bit} \approx 2.46$~Gbps.

\smallskip
\noindent\textbf{[A2] BLE share of the raw rate.} The BLE ceiling as a fraction of that raw
rate: $4\text{~Mbps} / 2.46\text{~Gbps} \approx 0.16\%$.

\smallskip
\noindent\textbf{[A3] Minimum spike-band sample rate.} By the Nyquist--Shannon sampling
theorem, faithfully capturing a band reaching $\sim$6~kHz requires
$f_s \geq 2 \times 6\text{~kHz} = 12$~kS/s.

\smallskip
\noindent\textbf{[A4] Raw snippet size.} A 40-sample spike snippet at 16-bit resolution is
$40 \times 16\text{~bit} = 640\text{~bit} = 80$~bytes, against the five-byte
identity-plus-timestamp packet of T11.

% ======= References ======= %
\phantomsection
\addcontentsline{toc}{section}{References}
\begin{thebibliography}{24}

\bibitem{ref1} E. R. Kandel, J. D. Koester, S. H. Mack, and S. A. Siegelbaum, \emph{Principles of Neural Science}, 6th ed. New York, NY, USA: McGraw-Hill, 2021.

\bibitem{ref2} E.-A. Seyfarth, ``Julius Bernstein (1839--1917): Pioneer neurobiologist and biophysicist,'' \emph{Biol. Cybern.}, vol. 94, no. 1, pp. 2--8, Jan. 2006.

\bibitem{ref3} J. J. Jun \emph{et al.}, ``Fully integrated silicon probes for high-density recording of neural activity,'' \emph{Nature}, vol. 551, no. 7679, pp. 232--236, Nov. 2017.

\bibitem{ref4} N. A. Steinmetz \emph{et al.}, ``Neuropixels 2.0: A miniaturized high-density probe for stable, long-term brain recordings,'' \emph{Science}, vol. 372, no. 6539, Art. no. eabf4588, Apr. 2021.

\bibitem{ref5} N. A. Steinmetz, P. Zatka-Haas, M. Carandini, and K. D. Harris, ``Distributed coding of choice, action and engagement across the mouse brain,'' \emph{Nature}, vol. 576, no. 7786, pp. 266--273, Dec. 2019.

\bibitem{ref6} IMEC and the Neuropixels Consortium, ``Neuropixels: Next-generation high-density neural probes,'' Leuven, Belgium, 2024. [Online]. Available: https://www.neuropixels.org

\bibitem{ref7} Intan Technologies, ``RHD2000-series digital electrophysiology interface chips (RHD2216/RHD2132),'' Datasheet, Los Angeles, CA, USA, 2021.

\bibitem{ref8} Intan Technologies, ``RHD2164 64-channel digital electrophysiology interface chip,'' Datasheet, Los Angeles, CA, USA, 2021.

\bibitem{ref9} Bluetooth Special Interest Group, ``Bluetooth core specification, version 6.0,'' Kirkland, WA, USA, 2024.

\bibitem{ref10} J. Tosi, F. Taffoni, M. Santacatterina, R. Sannino, and D. Formica, ``Performance evaluation of Bluetooth Low Energy: A systematic review,'' \emph{Sensors}, vol. 17, no. 12, Art. no. 2898, Dec. 2017.

\bibitem{ref11} H. G. Rey, C. Pedreira, and R. Quian Quiroga, ``Past, present and future of spike sorting techniques,'' \emph{Brain Res. Bull.}, vol. 119, pp. 106--117, Oct. 2015.

\bibitem{ref12} M. S. Lewicki, ``A review of methods for spike sorting: The detection and classification of neural action potentials,'' \emph{Network: Comput. Neural Syst.}, vol. 9, no. 4, pp. R53--R78, Nov. 1998.

\bibitem{ref13} N. Tasneem \emph{et al.}, ``A duty-cycled wireless optogenetic headstage for freely-moving rodents,'' \emph{IEEE Sensors J.}, vol. 23, no. 21, Nov. 2023.

\bibitem{ref14} G. Gagnon-Turcotte, Y. LeChasseur, C. Bories, Y. Messaddeq, Y. De Koninck, and B. Gosselin, ``A wireless optogenetic headstage with multichannel electrophysiological recording and in-situ spike detection and compression,'' \emph{IEEE Trans. Biomed. Circuits Syst.}, vol. 11, no. 1, pp. 129--141, Feb. 2017.

\bibitem{ref15} Z. Chen \emph{et al.}, ``A 128-channel spike-sorting SoC with level-crossing digitization and spatiotemporal feature extraction,'' \emph{IEEE Trans. Biomed. Circuits Syst.}, 2026.

\bibitem{ref16} Z. Jiang \emph{et al.}, ``A 1024-channel 583\,nW/channel spike-sorting SoC for high-density neural recording,'' \emph{IEEE Trans. Circuits Syst. I}, vol. 73, no. 8, 2026.

\bibitem{ref17} R. Lukito \emph{et al.}, ``A spike-sorting SoC with analog computing-in-memory autoencoder feature extraction,'' \emph{IEEE J. Solid-State Circuits}, vol. 61, no. 4, 2026.

\bibitem{ref18} S. Luan \emph{et al.}, ``Compact standalone platform for neural recording with real-time spike sorting and data logging,'' \emph{J. Neural Eng.}, vol. 15, no. 4, Art. no. 046014, 2018.

\bibitem{ref19} Z. Wu \emph{et al.}, ``A real-time unsupervised OSort spike-sorting processor implemented on FPGA,'' \emph{J. Neural Eng.}, vol. 20, no. 6, Art. no. 066015, 2023.

\bibitem{ref20} J. Navajas, D. Y. Barsakcioglu, A. Eftekhar, A. Jackson, T. G. Constandinou, and R. Quian Quiroga, ``Minimum requirements for accurate and efficient real-time on-chip spike sorting,'' \emph{J. Neurosci. Methods}, vol. 230, pp. 51--64, Jun. 2014.

\bibitem{ref21} R. R. Harrison, ``The design of integrated circuits to observe brain activity,'' \emph{Proc. IEEE}, vol. 96, no. 7, pp. 1203--1216, Jul. 2008.

\bibitem{ref22} G. Buzs\'aki, C. A. Anastassiou, and C. Koch, ``The origin of extracellular fields and currents---EEG, ECoG, LFP and spikes,'' \emph{Nat. Rev. Neurosci.}, vol. 13, no. 6, pp. 407--420, Jun. 2012.

\bibitem{ref23} Nordic Semiconductor, ``nRF54LM20B objective product specification,'' Trondheim, Norway, 2024.

\bibitem{ref24} M. Pachitariu, S. Sridhar, J. Pennington, and C. Stringer, ``Spike sorting with Kilosort4,'' \emph{Nat. Methods}, vol. 21, no. 5, pp. 914--921, May 2024.

\end{thebibliography}

\end{document}
